\documentclass[12pt,a4paper]{article}
\usepackage[UTF8]{ctex}
\usepackage{amsmath,amssymb,amsthm}
\usepackage{geometry}
\usepackage{bm}
\usepackage{hyperref}

\geometry{margin=2.5cm}

\title{空白区域与Wrench锥的关系：为什么力矩符号不定仍然不行？}
\author{现代机器人学：第12章 抓取与操作\\
详细解释}
\date{}

\begin{document}

\maketitle

\section{问题的核心}

\subsection{看似矛盾的现象}

\begin{itemize}
    \item \textbf{空白区域}：表示力矩符号取决于系数选择，可以产生任意符号的力矩
    \item \textbf{但}：wrench $(0, 0, mg)$不在复合wrench锥中，就不可行
    \item \textbf{困惑}：如果空白区域表示"灵活性大"，为什么wrench不在锥中就不行？
\end{itemize}

\subsection{关键误解}

\textbf{误解}：空白区域表示可以产生"任意"wrench。

\textbf{正确理解}：空白区域只表示对于\textbf{平面中的某个点}，可以产生\textbf{任意符号的力矩}，但这不意味着可以产生\textbf{任意的wrench}。

\section{Wrench vs 力矩}

\subsection{Wrench的完整定义}

平面wrench是一个\textbf{三维向量}：
\begin{equation}
F = (m_z, f_x, f_y)
\end{equation}

其中：
\begin{itemize}
    \item $m_z$：关于$\hat{z}$轴的力矩
    \item $f_x$：$\hat{x}$方向的力
    \item $f_y$：$\hat{y}$方向的力
\end{itemize}

\subsection{力矩标记只关注力矩}

\textbf{关键点}：力矩标记方法只关注wrench对\textbf{平面中某个点}产生的\textbf{力矩}，而不考虑wrench的\textbf{力分量}。

\subsubsection{力矩的计算}

对于平面中的点$p = (x, y)$，wrench $F = (m_z, f_x, f_y)$对该点产生的力矩为：
\begin{equation}
m_z(p) = m_z + x f_y - y f_x
\end{equation}

\subsubsection{空白区域的含义}

空白区域表示：对于点$p$，正线性组合可以产生\textbf{任意符号的力矩}$m_z(p)$。

但这\textbf{不意味着}可以产生\textbf{任意的wrench} $F = (m_z, f_x, f_y)$。

\section{为什么空白区域不能保证Wrench在锥中？}

\subsection{三个约束 vs 一个自由度}

\subsubsection{Wrench的三个分量}

要产生特定的wrench $F = (m_z, f_x, f_y)$，需要同时满足：
\begin{enumerate}
    \item 力矩分量：$m_z = m_{z,\text{desired}}$
    \item x方向力：$f_x = f_{x,\text{desired}}$
    \item y方向力：$f_y = f_{y,\text{desired}}$
\end{enumerate}

\subsubsection{空白区域只提供一个自由度}

空白区域只告诉我们：
\begin{itemize}
    \item 对于点$p$，可以产生\textbf{任意符号的力矩}$m_z(p)$
    \item 这是一个\textbf{自由度}
    \item 但wrench需要\textbf{三个分量}都正确
\end{itemize}

\subsection{具体例子：例12.9}

\subsubsection{所需的Wrench}

要使块保持静止，需要：
\begin{equation}
F_{\text{required}} = (0, 0, mg)
\end{equation}

即：
\begin{itemize}
    \item 力矩：$m_z = 0$
    \item x方向力：$f_x = 0$
    \item y方向力：$f_y = mg$（向上）
\end{itemize}

\subsubsection{复合Wrench锥的限制}

复合wrench锥$\mathcal{WC} = \text{pos}(\{F_1, F_2, F_3\})$，其中：
\begin{align}
F_1 &= \frac{1}{\sqrt{2}}(-4, -1, 1) \\
F_2 &= \frac{1}{\sqrt{2}}(-2, 1, 1) \\
F_3 &= (1, -1, 0)
\end{align}

\subsubsection{检查是否在锥中}

需要检查是否存在$k_1, k_2, k_3 \geq 0$，使得：
\begin{equation}
k_1 F_1 + k_2 F_2 + k_3 F_3 = (0, 0, mg)
\end{equation}

展开：
\begin{align}
-\frac{4k_1}{\sqrt{2}} - \frac{2k_2}{\sqrt{2}} + k_3 &= 0 \quad \text{（力矩）} \\
-\frac{k_1}{\sqrt{2}} + \frac{k_2}{\sqrt{2}} - k_3 &= 0 \quad \text{（x方向力）} \\
\frac{k_1}{\sqrt{2}} + \frac{k_2}{\sqrt{2}} &= mg \quad \text{（y方向力）}
\end{align}

\subsubsection{求解约束}

从第一个和第二个方程：
\begin{align}
-\frac{4k_1}{\sqrt{2}} - \frac{2k_2}{\sqrt{2}} + k_3 &= 0 \\
-\frac{k_1}{\sqrt{2}} + \frac{k_2}{\sqrt{2}} - k_3 &= 0
\end{align}

相加：
\begin{equation}
-\frac{5k_1}{\sqrt{2}} - \frac{k_2}{\sqrt{2}} = 0
\end{equation}

由于$k_1, k_2 \geq 0$，这要求$k_1 = k_2 = 0$。

但从第三个方程：
\begin{equation}
\frac{k_1 + k_2}{\sqrt{2}} = mg > 0
\end{equation}

这要求$k_1 + k_2 > 0$，与$k_1 = k_2 = 0$矛盾。

\textbf{因此}：不存在非负系数使得$k_1 F_1 + k_2 F_2 + k_3 F_3 = (0, 0, mg)$。

\subsubsection{空白区域的作用}

即使对于某些点$p$，wrench锥产生空白区域（可以产生任意符号的力矩$m_z(p)$），这\textbf{不意味着}可以产生wrench $(0, 0, mg)$。

因为：
\begin{itemize}
    \item 空白区域只保证\textbf{力矩}$m_z(p)$的灵活性
    \item 但wrench还需要\textbf{力分量}$f_x$和$f_y$正确
    \item 这三个约束可能无法同时满足
\end{itemize}

\section{力矩标记的局限性}

\subsection{力矩标记只提供部分信息}

力矩标记方法：
\begin{itemize}
    \item \textbf{提供}：对于每个点$p$，力矩$m_z(p)$的符号信息
    \item \textbf{不提供}：wrench的力分量$f_x$和$f_y$的信息
    \item \textbf{不提供}：wrench是否在锥中的完整信息
\end{itemize}

\subsection{完整判断需要什么？}

要判断wrench $F$是否在复合锥$\mathcal{WC}$中，需要：

\begin{enumerate}
    \item \textbf{检查线性组合}：是否存在$k_i \geq 0$使得$\sum_i k_i F_i = F$
    \item \textbf{或者检查锥的边界}：$F$是否在锥的边界或内部
    \item \textbf{或者使用力矩标记}：检查$F$的箭头是否在标记区域内（但这只提供部分信息）
\end{enumerate}

\section{为什么空白区域仍然有用？}

\subsection{力闭合的判断}

虽然空白区域不能直接判断wrench是否在锥中，但它对\textbf{力闭合}的判断很有用：

\begin{itemize}
    \item \textbf{力闭合要求}：没有一致标记的区域
    \item \textbf{原因}：如果所有区域都是空白（不一致），则wrench锥可以产生任意方向的wrench
    \item \textbf{但}：这仍然不意味着可以产生\textbf{任意大小}的wrench
\end{itemize}

\subsection{力矩标记的适用场景}

力矩标记适用于：
\begin{itemize}
    \item \textbf{判断力闭合}：检查是否有一致标记区域
    \item \textbf{可视化wrench锥}：理解wrench锥的几何结构
    \item \textbf{定性分析}：理解哪些方向的wrench可能产生
\end{itemize}

但不适用于：
\begin{itemize}
    \item \textbf{精确判断}：判断特定wrench是否在锥中
    \item \textbf{定量分析}：计算所需的接触力
\end{itemize}

\section{正确的理解}

\subsection{空白区域的含义}

空白区域表示：
\begin{itemize}
    \item 对于点$p$，可以产生\textbf{任意符号的力矩}$m_z(p)$
    \item 这是一个\textbf{自由度}
    \item 但这\textbf{不意味着}可以产生\textbf{任意的wrench}
\end{itemize}

\subsection{Wrench在锥中的条件}

Wrench $F$在复合锥$\mathcal{WC}$中的条件是：
\begin{equation}
F \in \mathcal{WC} = \text{pos}(\{F_1, F_2, \ldots, F_n\})
\end{equation}

即存在$k_i \geq 0$使得：
\begin{equation}
F = \sum_{i=1}^n k_i F_i
\end{equation}

这需要\textbf{三个分量}都同时满足，而不仅仅是力矩。

\subsection{例12.9的情况}

\begin{itemize}
    \item 即使某些区域是空白（力矩灵活）
    \item 但wrench $(0, 0, mg)$的三个分量无法同时满足
    \item 因此不在复合锥中
    \item 静态抓取不可行
\end{itemize}

\section{关键要点总结}

\begin{enumerate}
    \item \textbf{空白区域}：只表示力矩$m_z(p)$的灵活性，不表示整个wrench的灵活性
    
    \item \textbf{Wrench是三维的}：需要同时满足$m_z$、$f_x$、$f_y$三个分量
    
    \item \textbf{力矩标记的局限性}：只提供力矩信息，不提供力的信息
    
    \item \textbf{判断wrench在锥中}：需要检查三个分量是否都能通过正线性组合实现
    
    \item \textbf{例12.9}：即使有空白区域，wrench $(0, 0, mg)$的三个分量无法同时满足，因此不在锥中
    
    \item \textbf{力矩标记的用途}：主要用于力闭合判断和可视化，不用于精确判断特定wrench
\end{enumerate}

\section{实际应用}

理解这个区别有助于：

\begin{itemize}
    \item \textbf{正确使用力矩标记}：理解其适用场景和局限性
    \item \textbf{精确分析}：使用代数方法判断特定wrench是否可行
    \item \textbf{设计抓取}：理解需要同时考虑力矩和力的约束
    \item \textbf{避免误解}：不要将力矩的灵活性与wrench的灵活性混淆
\end{itemize}

\end{document}

